\documentclass[12pt,a4paper]{report}

\usepackage[T1]{fontenc}
\usepackage[utf8]{inputenc}
\usepackage[serbian]{babel}

\usepackage{amsmath,amssymb,amsthm}
\usepackage{graphicx}
\usepackage{geometry}
\usepackage{setspace}
\usepackage{titlesec}
\usepackage{tocloft}

\geometry{margin=2.5cm}
\onehalfspacing

\begin{document}

\begin{titlepage}
\centering

{\Large Univerzitet Singidunum \par}
\vspace{0.5cm}

\vspace{3cm}

{\Huge \textbf{ODREĐIVANJE n-TOG STEPENA CIRKULARNIH MATRICA ČIJI SU ELEMENTI POZNATI NIZOVI} \par}

\vspace{3cm}

\begin{flushleft}
\large
Mentor: prof. dr Biljana Radičić

\vspace{1cm}

Kandidat: Marko Milošević
\end{flushleft}

\vfill

{\large Beograd, 2026.}

\end{titlepage}

\tableofcontents
\newpage

\chapter{Uvod}

Cirkularne matrice predstavljaju posebnu klasu matrica koje imaju značajnu primenu u linearnoj algebri, teoriji signala, obradi slike, numeričkoj matematici i teoriji diskretnih sistema. Njihova osnovna karakteristika jeste da je svaki naredni red dobijen cikličnim pomeranjem prethodnog reda za jedno mesto.

Poseban značaj cirkularnih matrica proizilazi iz činjenice da se njihove sopstvene vrednosti i sopstveni vektori mogu eksplicitno odrediti korišćenjem korena jedinice i Furijeove transformacije.

U ovom radu razmatraće se određivanje n-tog stepena cirkularnih matrica čiji su elementi definisani pomoću poznatih nizova:

\begin{itemize}
    \item aritmetičkog niza,
    \item geometrijskog niza,
    \item Fibonačijevog niza.
\end{itemize}

Cilj rada jeste izvođenje opštih formula za određivanje stepena cirkularnih matrica i analiza njihovog ponašanja.

\chapter{Osnovni pojmovi i definicije}

\section{Matrice}

Matrica reda $m \times n$ predstavlja pravougaoni raspored elemenata:

\[
A = [a_{ij}]_{m \times n}.
\]

Kvadratna matrica reda $n$ ima jednak broj vrsta i kolona.

\section{Stepen matrice}

Za kvadratnu matricu $A$ definiše se:

\[
A^n = \underbrace{A \cdot A \cdot \dots \cdot A}_{n\text{ puta}}.
\]

Za $n=0$ važi:

\[
A^0 = I,
\]

gde je $I$ jedinična matrica.

\section{Poznati nizovi}

\subsection{Aritmetički niz}

Aritmetički niz definisan je formulom:

\[
a_n = a_1 + (n-1)d.
\]

\subsection{Geometrijski niz}

Geometrijski niz definisan je formulom:

\[
a_n = a_1 q^{n-1}.
\]

\subsection{Fibonačijev niz}

Fibonačijev niz definiše se rekurentno:

\[
F_n = F_{n-1} + F_{n-2},
\]

uz početne uslove:

\[
F_0 = 0, \qquad F_1 = 1.
\]

\chapter{Cirkularne matrice}

\section{Definicija}

Cirkularna matrica reda $n$ definiše se kao:

\[
C = \operatorname{circ}(c_0,c_1,c_2,\dots,c_{n-1}).
\]

Opšti oblik cirkularne matrice je:

\[
C=
\begin{bmatrix}
c_0 & c_1 & c_2 & \cdots & c_{n-1}\\
c_{n-1} & c_0 & c_1 & \cdots & c_{n-2}\\
\vdots & \vdots & \vdots & \ddots & \vdots\\
c_1 & c_2 & c_3 & \cdots & c_0
\end{bmatrix}.
\]

\section{Osnovna svojstva}

Važe sledeća svojstva:

\begin{enumerate}
    \item zbir dve cirkularne matrice je cirkularna matrica,
    \item proizvod dve cirkularne matrice je cirkularna matrica,
    \item inverzna matrica cirkularne matrice je takođe cirkularna,
    \item cirkularne matrice međusobno komutiraju.
\end{enumerate}

\section{Veza sa Furijeovom transformacijom}

Cirkularne matrice dijagonalizuju se pomoću Furijeove matrice:

\[
C = F^{-1}DF,
\]

gde je:

\begin{itemize}
    \item $F$ Furijeova matrica,
    \item $D$ dijagonalna matrica sopstvenih vrednosti.
\end{itemize}

\chapter{Sopstvene vrednosti i sopstveni vektori}

Neka je:

\[
C = \operatorname{circ}(c_0,c_1,\dots,c_{n-1}).
\]

Sopstvene vrednosti matrice određuju se formulom:

\[
\lambda_k = \sum_{j=0}^{n-1} c_j \omega_n^{jk},
\]

gde je:

\[
\omega_n = e^{\frac{2\pi i}{n}}.
\]

Sopstveni vektori imaju oblik:

\[
v_k = (1,\omega_n^k,\omega_n^{2k},\dots,\omega_n^{(n-1)k}).
\]

Pošto su cirkularne matrice dijagonalizabilne:

\[
C^m = F^{-1}D^mF.
\]

\chapter{Određivanje n-tog stepena cirkularnih matrica}

Neka je:

\[
C = F^{-1}DF.
\]

Tada važi:

\[
C^n = F^{-1}D^nF.
\]

Pošto je:

\[
D = \operatorname{diag}(\lambda_0,\lambda_1,\dots,\lambda_{n-1}),
\]

sledi:

\[
D^n = \operatorname{diag}(\lambda_0^n,\lambda_1^n,\dots,\lambda_{n-1}^n).
\]

Dakle, određivanje n-tog stepena matrice svodi se na stepenovanje sopstvenih vrednosti.

\chapter{Cirkularne matrice sa elementima aritmetičkog niza}

Razmatrajmo matricu:

\[
A = \operatorname{circ}(a,a+d,a+2d).
\]

Njena forma je:

\[
A=
\begin{bmatrix}
a & a+d & a+2d \\
a+2d & a & a+d \\
a+d & a+2d & a
\end{bmatrix}.
\]

\section{Sopstvene vrednosti}

Dobijamo:

\[
\lambda_k = a + (a+d)\omega^k + (a+2d)\omega^{2k}.
\]

Za $k=0$:

\[
\lambda_0 = 3a+3d.
\]

\section{Primer}

Neka je:

\[
a=1, \qquad d=1.
\]

Tada je:

\[
A = \operatorname{circ}(1,2,3).
\]

Odnosno:

\[
A=
\begin{bmatrix}
1 & 2 & 3 \\
3 & 1 & 2 \\
2 & 3 & 1
\end{bmatrix}.
\]

Računamo:

\[
A^2=
\begin{bmatrix}
13 & 13 & 10 \\
10 & 13 & 13 \\
13 & 10 & 13
\end{bmatrix}.
\]

\chapter{Cirkularne matrice sa elementima geometrijskog niza}

Posmatrajmo matricu:

\[
G = \operatorname{circ}(a,aq,aq^2).
\]

Odnosno:

\[
G=
\begin{bmatrix}
a & aq & aq^2 \\
aq^2 & a & aq \\
aq & aq^2 & a
\end{bmatrix}.
\]

Sopstvene vrednosti su:

\[
\lambda_k = a(1+q\omega^k+q^2\omega^{2k}).
\]

\section{Primer}

Neka je:

\[
a=1, \qquad q=2.
\]

Tada je:

\[
G = \operatorname{circ}(1,2,4).
\]

Kvadrat matrice iznosi:

\[
G^2=
\begin{bmatrix}
17 & 20 & 16 \\
16 & 17 & 20 \\
20 & 16 & 17
\end{bmatrix}.
\]

\chapter{Cirkularne matrice sa elementima Fibonačijevog niza}

Razmatrajmo matricu reda 4:

\[
F = \operatorname{circ}(F_1,F_2,F_3,F_4).
\]

Pošto su:

\[
F_1=1, \quad F_2=1, \quad F_3=2, \quad F_4=3,
\]

dobijamo:

\[
F = \operatorname{circ}(1,1,2,3).
\]

Odnosno:

\[
F=
\begin{bmatrix}
1 & 1 & 2 & 3 \\
3 & 1 & 1 & 2 \\
2 & 3 & 1 & 1 \\
1 & 2 & 3 & 1
\end{bmatrix}.
\]

\section{Sopstvene vrednosti}

Koreni jedinice četvrtog reda su:

\[
1,-1,i,-i.
\]

Dobijamo:

\[
\lambda_0 = 7,
\]

\[
\lambda_1 = -1,
\]

\[
\lambda_2 = -1+i,
\]

\[
\lambda_3 = -1-i.
\]

\section{Kvadrat matrice}

\[
F^2=
\begin{bmatrix}
14 & 14 & 15 & 11 \\
11 & 14 & 14 & 15 \\
15 & 11 & 14 & 14 \\
14 & 15 & 11 & 14
\end{bmatrix}.
\]

\chapter{Primeri matrica reda 3 i 4}

\section{Primer matrice reda 3}

Neka je:

\[
B = \operatorname{circ}(2,3,4).
\]

Tada je:

\[
B=
\begin{bmatrix}
2 & 3 & 4 \\
4 & 2 & 3 \\
3 & 4 & 2
\end{bmatrix}.
\]

Računamo:

\[
B^2=
\begin{bmatrix}
28 & 28 & 25 \\
25 & 28 & 28 \\
28 & 25 & 28
\end{bmatrix}.
\]

\section{Primer matrice reda 4}

Neka je:

\[
C = \operatorname{circ}(1,2,3,4).
\]

Odnosno:

\[
C=
\begin{bmatrix}
1 & 2 & 3 & 4 \\
4 & 1 & 2 & 3 \\
3 & 4 & 1 & 2 \\
2 & 3 & 4 & 1
\end{bmatrix}.
\]

Kvadrat matrice:

\[
C^2=
\begin{bmatrix}
30 & 24 & 22 & 24 \\
24 & 30 & 24 & 22 \\
22 & 24 & 30 & 24 \\
24 & 22 & 24 & 30
\end{bmatrix}.
\]

\chapter{Primene cirkularnih matrica}

Cirkularne matrice imaju značajnu primenu u:

\begin{itemize}
    \item obradi digitalnih signala,
    \item teoriji kodovanja,
    \item numeričkoj analizi,
    \item obradi slike,
    \item kriptografiji,
    \item rešavanju diferencnih jednačina.
\end{itemize}

Njihova glavna prednost jeste mogućnost efikasnog računanja pomoću Furijeove transformacije.

\chapter{Zaključak}

U radu su analizirane cirkularne matrice čiji su elementi definisani pomoću poznatih nizova. Pokazano je da se određivanje n-tog stepena matrice može efikasno izvršiti korišćenjem dijagonalizacije.

Razmatrani su slučajevi:

\begin{itemize}
    \item aritmetičkog niza,
    \item geometrijskog niza,
    \item Fibonačijevog niza.
\end{itemize}

Analizirani primeri reda 3 i 4 pokazali su da se struktura cirkularnih matrica zadržava i nakon stepenovanja.

Dalja istraživanja mogu obuhvatiti:

\begin{itemize}
    \item matrice višeg reda,
    \item druge poznate nizove,
    \item numeričke algoritme za brzo računanje stepena.
\end{itemize}

\chapter{Literatura}

\begin{enumerate}
    \item Davis, P. J., \textit{Circulant Matrices}, Wiley, New York, 1979.

    \item Horn, R. A., Johnson, C. R., \textit{Matrix Analysis}, Cambridge University Press, 2012.

    \item Gray, R. M., \textit{Toeplitz and Circulant Matrices: A Review}, Now Publishers Inc., 2006.

    \item Strang, G., \textit{Linear Algebra and Its Applications}, Brooks/Cole, 2006.

    \item Koshy, T., \textit{Fibonacci and Lucas Numbers with Applications}, Wiley, 2001.

    \item Meyer, C. D., \textit{Matrix Analysis and Applied Linear Algebra}, SIAM, 2000.
\end{enumerate}

\end{document}
